\documentclass[11pt,a4paper]{article}
\usepackage[margin=25mm]{geometry}
\usepackage{amsmath,amssymb,amsthm,booktabs,hyperref}
\newtheorem{lemma}{Lemma}\newtheorem{proposition}[lemma]{Proposition}
\newcommand{\PGL}{\operatorname{PGL}}
\newcommand{\SL}{\operatorname{SL}}
\title{Exhausting the cubic matrix equations}
\author{Human Compression project: L023 working proof}\date{}
\begin{document}\maketitle
This unreviewed component recovers, with explicit corrections, the polynomial
construction and finite table in Leemans--Stokes, arXiv:2208.08215v2,
Section~4. It does not settle the separate rank-four incidence assertions.
The equations and existence construction belong to the source authors.
The orbit exhaustion and the characteristic-two correction below are the
present argument. No conjecture, whole-original, formalization or
compression factor is claimed.

Put $K=\mathbb F_{q^3}$, $k=\mathbb F_q$, and let $\tau$ raise entries
to their $q$th powers. Let $J=\left(\begin{smallmatrix}0&1\\-1&0\end{smallmatrix}\right)$.
The first three source equations are equivalent to
\[
 R=\begin{pmatrix}a&b\\c&-a\end{pmatrix},\quad
 \det R=1,\quad R^{\tau^2}=RR^\tau.                 \tag{1}
\]
Here superscripts $\tau$ denote the field operation, not matrix powers.
The source's trace equation is redundant in (1).

\begin{lemma}[Odd characteristic]
If $q$ is odd and (1) has a solution, all its solutions form one free
orbit under conjugation by $\PGL_2(k)$. There are exactly $q(q^2-1)$
solutions.
\end{lemma}
\begin{proof}
Write $S=R^\tau$ and $T=R^{\tau^2}=RS$. Cayley--Hamilton gives
$R^2=S^2=T^2=-I$. Consequently $SR=-RS$. The matrices
$I,R,S,RS$ form a basis of $M_2(K)$: under the trace pairing their
Gram matrix is diagonal with entries $2,-2,-2,-2$, all nonzero.
Their multiplication is determined by $R^2=S^2=-I$ and $SR=-RS$.
For any second solution $R'$, the correspondence taking this basis to
$I,R',S',R'S'$ is therefore a $K$-algebra automorphism of $M_2(K)$.

Every such algebra automorphism is conjugation by a matrix $C\in\mathrm{GL}_2(K)$.
For completeness, choose a basis adapted to the images of the two
orthogonal rank-one idempotents $E_{11},E_{22}$. In that basis the
images of $E_{12},E_{21}$ are respectively $uE_{12},u^{-1}E_{21}$,
by their products with the idempotents and each other. Rescaling one
basis vector gives the claimed conjugation.
Thus $R'=CRC^{-1}$ and $S'=CSC^{-1}$. Apply $\tau$ to these
identities and compare with the identities for $S',T'$. The matrix
$C^{-1}C^\tau$ commutes with $S,T$, hence with all of $M_2(K)$,
and is scalar. Divide $C$ by any one nonzero entry. All its entry
ratios are fixed by $\tau$, so this normalized matrix lies in
$\mathrm{GL}_2(k)$. This proves transitivity under $\PGL_2(k)$.
A projectivity stabilizing $R$ also stabilizes $S$, hence is scalar
by the same basis calculation. The action is free, giving the count.
\end{proof}

\begin{lemma}[Characteristic two]
If $q$ is even, the nonscalar solutions of (1) are parametrized without
repetition by a point of $\mathbb P^1(k)$ and a nonzero element of
$\ker\operatorname{Tr}_{K/k}$. Their number is $(q+1)(q^2-1)$.
The only additional solution is $I$.
\end{lemma}
\begin{proof}
A nonscalar trace-zero determinant-one matrix is a nonidentity unipotent
involution, with a unique fixed point on $\mathbb P^1(K)$.
With $S=R^\tau$ the identity $(RS)^2=I$ says that $R,S$ commute.
They are distinct, since $R=S$ would give $R^{\tau^2}=I$.
Two commuting nonidentity unipotents in $\SL_2(K)$ have the same
fixed point: after conjugating one to upper triangular form, its
centralizer is the group of upper unitriangular matrices. Their shared
fixed point is preserved by $\tau$, hence belongs to $\mathbb P^1(k)$.

Choose, once for each such point $P$, a matrix $C_P\in\mathrm{GL}_2(k)$
carrying infinity to $P$. Every solution with this point has the unique
form
\[
 R=C_P\begin{pmatrix}1&t\\0&1\end{pmatrix}C_P^{-1},
 \qquad t\ne0.
\]
Equation (1) is now exactly $t^{q^2}=t+t^q$, the trace-zero condition.
The trace map is a nonzero $k$-linear map from a three-dimensional
space to $k$ (on $k$ it is the identity), so its kernel has $q^2$
elements. This gives all the stated parameters without repetition.
The scalar case forces $R=I$, which does satisfy (1).
\end{proof}

\begin{proposition}[Primitive eigenvalues]
Let $Q=q^3\geq8$. Suppose $R$ satisfies (1) and $RJ$ has order
$Q-1$ or $Q+1$ in $\SL_2(K)$. Its projective images, together with
$R^\tau$, generate $\operatorname{PSL}_2(K)$ and give a Wilson
class~III map. The three periods are
\[
                \frac{Q-1}{\gcd(2,Q-1)}
 \quad\hbox{or}\quad
                \frac{Q+1}{\gcd(2,Q-1)},                \tag{2}
\]
respectively.
\end{proposition}
\begin{proof}
The projective images $r,s$ of $R,R^\tau$ are distinct commuting
involutions, and $\tau$ cycles $r,s,rs$ while fixing the image $j$
of $J$. Let $d=\gcd(2,Q-1)$ and $m=(Q\pm1)/d$. The order
assumption makes $\langle r,j\rangle$ the full dihedral torus
normalizer of order $2m$. By the retained Dickson classification,
this is maximal for the cubic field sizes $Q\geq8$ under consideration.
Thus $L=\langle r,j,s\rangle$ is either that dihedral group or the
whole projective special linear group.

The former alternative is impossible. If $m$ is odd the dihedral group
has no Klein four subgroup. If $m$ is even, each Klein four subgroup
contains its unique central involution. But $L$ is $\tau$-stable, and
$\tau$ cycles the three involutions of $\langle r,s\rangle$, so
cannot fix that characteristic central involution. This proves generation.
In odd characteristic the projective order is half the indicated matrix
order; in even characteristic it is unchanged, proving (2).
Finally $\tau$ is a noninner field automorphism. A Wilson duality would
combine with it to give a faithful $S_3$ of marked automorphisms; its
commutator would die in the abelian outer automorphism group, contradicting
the nontrivial field image of $\tau$. Thus the map has class~III.
\end{proof}
The statement uses a \emph{primitive} $(Q\pm1)$st root of unity in the
source's fourth equation. Allowing roots of smaller order would not give
the required matrix order. Indeed a determinant-one matrix has characteristic
polynomial $x^2-\operatorname{tr}(RJ)x+1$; the trace $\xi+\xi^{-1}$
for a primitive root gives the two distinct eigenvalues $\xi,\xi^{-1}$.
The printed ambient $\operatorname{PSL}_2(q)$ in the proof of source
Lemma~5 must be $\operatorname{PSL}_2(q^3)$. The half-orders printed after
its table apply only in odd characteristic.

\section*{The complete finite comparison}
The whole required program \texttt{polynomial\_table.g} works in GAP's
exact finite fields. For odd $q$ it finds one literal solution by scanning
$a,b$ and setting $c=(-1-a^2)/b$, with $b\ne0$. That scan is used only
to find a seed; it need not exhaust all possible seeds. The program checks
that a seed is found for every odd $q$ in the table. It then conjugates
this seed by the following complete list of projectivities over $k$:
\[
 \begin{pmatrix}x&y\\0&1\end{pmatrix}\ (x\ne0),\qquad
 \begin{pmatrix}x&y\\1&z\end{pmatrix}\ (xz\ne y).
\]
Their respective sizes are $q(q-1)$ and $q^2(q-1)$. Normalizing the
first nonzero entry of the bottom row proves both completeness and
uniqueness. The odd-characteristic lemma therefore proves exhaustion.
For even $q$ the program uses $C_\infty=I$ and
$C_x=\left(\begin{smallmatrix}x&1\\1&0\end{smallmatrix}\right)$
for $x\in k$, and every nonzero trace-zero parameter. The omitted scalar
solution $I$ has $IJ=J$ of order two, and contributes to neither column.

For each resulting matrix the program checks determinant one, trace zero,
and the full matrix equation (1), rejects duplicate matrices, checks the
proved total, and computes the exact order of $RJ$. The output records the
seed by exponents of GAP's $Z(q^3)$, with $-1$ for zero. Thus the finite
counts use the complete program and finite-field arithmetic as a
computational premise; they are not inferred from a sample or a rendered
table. No optional geometry probe is a premise.

\begin{center}
\begin{tabular}{rrrrrr}
\toprule
$q$&nonscalar solutions&$Q-1$&$Q+1$&printed $Q-1$&printed $Q+1$\\
\midrule
3&24&0&24&0&24\\
4&75&24&24&12&12\\
5&120&48&36&48&36\\
7&336&96&96&96&96\\
8&567&216&144&162&108\\
9&720&192&192&192&192\\
11&1320&324&216&324&216\\
13&2184&432&576&432&576\\
16&4335&960&1920&840&1680\\
17&4896&1440&624&1440&624\\
19&6840&1020&1440&1020&1440\\
23&12144&2952&2160&2952&2160\\
\bottomrule
\end{tabular}
\end{center}
In characteristic two, the excess counts are respectively $(12,12)$,
$(54,36)$ and $(120,240)$. In the complete census these are exactly
the qualifying solutions with $bc=0$. Subtracting those solutions
reproduces the three printed rows. The source does not impose $bc\ne0$;
these are therefore additional solutions to its printed equations under
the primitive-root interpretation. This observation identifies which
solutions account for the numerical discrepancy, without claiming to know
how the source authors' computation produced it.

\paragraph{Scope.}
This is a component of a prospective corrected recovery. The retained
inputs are the stated torus-normalizer part of Dickson's classification,
the projective automorphism description and Wilson's marked-map
correspondence. The two exhaustion lemmas are proved here. The displayed
finite values require the whole GAP program. No claim about full geometric
correlations of the rank-four construction follows from these calculations.
\end{document}
